\hspace{3mm}

\centerline{\bf \huge 8 класс}

\hspace{3mm}

\problem{1} Кирилл пришел на I ежегодную Математическую Олимпиаду ЭлМШ, на которой было предложено 5 задач. Кирилл очень торопился и мог потратить на Олимпиаду максимум час. Задачи, придуманные Расулом Владимировичем, Кирилл читает за 2 минуты, а решает за 37 минут. Задачи, придуманные Анджуром Аркадьевичем, он читает 2 минуты, а решает за 15 минут. Задачи, придуманные Очиром Владимировичем, он читает 10 минут и решает тоже за 10 минут. Очир Владимирович предложил 2 задачи, Анджур Аркадьевич одну, а Расул Владимирович две. Какое количество баллов он успел набрать, если задачи Очира Владимировича он мог отличить, даже не читая, а отличить задачи Анджура Аркадьевича от задач Расула Владимировича может только прочитав их. Хоть Кириллу при прочтении задачи Расула Владимировича попадались раньше задач Анджура Аркадьевича, но поступал он наилучшим способом. За задачи Очира Владимировича ученики получали по 4 балла, Анджура Аркадьевича - по 5 баллов, Расула Владимировича - по 8 баллов.

\begin{flushright}
    \scriptsize{$-$ Ты остался один. \\
                $-$ Одного нельзя приравнивать к нулю.}
\end{flushright}

\hspace{3mm}

\problem{2} Найдите цифры a и b, для которых $\sqrt{0,aaaaa\dotsc} = 0,bbbbb \dotsc$

\begin{flushright}
    \scriptsize{$-$ AAAAAAAAA \\
                Bomb Man}
\end{flushright}

\hspace{3mm}

\problem{3} В выпуклом четырёхугольнике $ELMS$ перпендикуляры, проведенные к серединам сторон $EL, LM$ и $SE$ пересекаются в точке $H$. Найдите периметр четырёхугольника $ELMS$, если $EM = 6$, а углы $\angle HES=\angle HLE=\angle HML=45^{\circ}$.

\begin{flushright}
    \scriptsize{Всем доброе утро!\\ Сегодня я что-то заблудился на жизненном пути...}
\end{flushright}



\problem{4}На трёх натуральных числах сидят по Шиншилле, Лососю и Киви. Каждый из этих зверьков умеет прыгать на какое-то натуральное расстояние вправо или влево. Первый на $a$, второй на $b$ и третий на $c$.

\newpage

\hspace{3mm}

\noindentКаждый из них начал прыгать влево до тех пор, пока может (то есть пока не встретится с ненатуральным числом), а затем вправо до бесконечности, причем, попадая на число, они ее окрашивают. Оказалось так, что все натуральные числа были покрашены зверьками, причем никакое не было покрашено дважды. Докажите, что $ \dfrac{1}{a} + \dfrac{1}{b} + \dfrac{1}{c} = 1$.

\begin{flushright}
    \scriptsize{— Зоро, куда нам идти?\\
                — Туда.\\
                — Отлично, идем в обратную сторону!}
\end{flushright}

\hspace{3mm}

\problem{5} Баир выписал на доску в порядке возрастания все рациональные дроби между нулём и единицей, знаменатели, которых не превосходят $n$. Пусть $\dfrac{a}{b}$ и $\dfrac{c}{d}$ - две соседние (несократимые) дроби, выписанные на доску. Докажите, что $|ad - bc| = 1$.

\begin{flushright}
    \scriptsize{$-$ Дурак тот, кто дерётся без шансов на победу.}
\end{flushright}


